\begin{derivation}[高斯波包的自由传播与方程核对]
在二阶包络模型中，连续动量谱给传播核
\begin{equation*}
 K(\zeta,\zeta';t)=\int\frac{\dd p}{2\pi\hbar}
 \exp\!\left[\frac{\ii p(\zeta-\zeta')}{\hbar}
             -\frac{\ii p^2t}{2m_\parallel\hbar}\right]
 =\sqrt{\frac{m_\parallel}{2\pi\ii\hbar t}}
 \exp\!\left[\frac{\ii m_\parallel(\zeta-\zeta')^2}{2\hbar t}\right],
 \quad t>0.
\end{equation*}
振荡高斯积分可先加正的收敛因子，再取其趋零极限。将初态代入 $\psi(\zeta,t)=\int K\psi(\zeta',0)\dd\zeta'$，整理为
\begin{equation*}
 -A\zeta'^2+B\zeta'+C,
 \quad A=a_0-\frac{\ii m_\parallel}{2\hbar t},\quad
 B=-\frac{\ii m_\parallel\zeta}{\hbar t},\quad
 C=\frac{\ii m_\parallel\zeta^2}{2\hbar t}.
\end{equation*}
$\Re A>0$ 保证积分收敛；使用 $\int\ee^{-Ax^2+Bx}\dd x=\sqrt{\pi/A}\ee^{B^2/(4A)}$，合并前因子与指数即得\eqnrefs{eq:electron-gaussian-tdse-solution}。

还可独立代回检验。写 $a(t)=a_0/d(t)$、$\mathcal N(t)=\mathcal N_0/\sqrt{d(t)}$，则
\begin{equation*}
 \dot a=-\frac{2\ii\hbar}{m_\parallel}a^2,\qquad
 \frac{\dot{\mathcal N}}{\mathcal N}=-\frac{\ii\hbar}{m_\parallel}a,
 \qquad
 \frac{\partial_\zeta^2\psi}{\psi}=-2a+4a^2\zeta^2.
\end{equation*}
含时 Schrödinger 方程左右两侧的常数项与 $\zeta^2$ 项于是分别相等。再用
$\Re(a_0/d)=\Re a_0/|d|^2$，对 $|\psi|^2$ 作高斯积分，归一化保持为一。所以传播后的函数同时满足初始条件、运动方程和归一化条件。
\end{derivation}
